\documentclass[11pt,letterpaper]{article}
\usepackage[margin=0.9in]{geometry}
\usepackage{amsmath,amssymb,amsthm,booktabs,microtype}
\usepackage[T1]{fontenc}
\usepackage{lmodern}
\usepackage[colorlinks=true,linkcolor=blue,citecolor=blue,urlcolor=blue]{hyperref}
\usepackage{fancyhdr}
\pagestyle{fancy}\fancyhf{}
\fancyhead[L]{\small Filtered MOND: a central tidal consistency relation}
\fancyhead[R]{\small Research draft}\fancyfoot[C]{\thepage}
\setlength{\headheight}{14pt}
\setlength{\parindent}{0pt}\setlength{\parskip}{5pt}
\newtheorem{proposition}{Proposition}
\newcommand{\ph}{\mathrm{ph}}\newcommand{\dd}{\mathrm{d}}
\newcommand{\ee}{\hat{\boldsymbol e}}\newcommand{\kk}{\hat{\boldsymbol k}}
\newcommand{\tr}{\operatorname{tr}}
\hypersetup{pdftitle={Why Banik and Zhao Were Right},pdfauthor={Carl Zimmerman},
pdfsubject={Conditional central tidal identity in action-derived filtered MOND}}
\begin{document}
\begin{center}
{\LARGE\bfseries Why Banik and Zhao Were Right\par}
\vspace{8pt}
{\large External-field structure and a central tidal\par
consistency relation in filtered MOND\par}
\vspace{12pt}
Carl Zimmerman\par
\small AI-assisted research draft; not peer reviewed\par
4 September 2026
\end{center}
\begin{abstract}
A modified-gravity law cannot be tested from an isolated spherical force curve
alone: its response to an ambient field also matters. We revisit the
external-field linearization of QUMOND and study an explicitly varied,
spatially filtered specialization of its two-potential action. The Newtonian
source is left unsmoothed; variation places an adjoint filter outside the
nonlinear phantom-density source. For a spherical source perturbing a constant,
nonzero constitutive background, isotropic filtering produces a finite central
phantom-potential Hessian whenever the stated spectral integral converges.
Its trace $D$ and axial quadrupole $Q_2$ obey a division-free consistency
identity independent of source mass and filter amplitude. For the inverse
partner of $\mu(y)=1-e^{-y}$, the ratio is
$Q_2/D=2y/\{5[3+3(y-1)e^{-y}-y]\}$ where $D\ne0$.
Symbolic identities and independent quadrature over 324 cases check this
result, with a deliberately anisotropic control violating it. The contribution
is a conditional central-response corollary, not a new MOND architecture,
an empirical detection, or a completed relativistic theory.
\end{abstract}
\section{What the title means}
Banik and Zhao derived the anisotropic external-field-dominated response of
AQUAL and QUMOND \cite{banik}. In QUMOND their equation (24) contains both
the value of the interpolation function and its logarithmic derivative.
Suppressing or retaining a scalar force enhancement therefore does not, by
itself, specify the directional response. This conditional mathematical point
is the sense in which the title says they were right; we do not adjudicate
their observational claims or any wide-binary controversy.

The parent two-potential action and external-field asymptotics are due to
Milgrom \cite{qumond}. His later generalization explicitly allows a possibly
nonlocal functional of the Newtonian potential and investigates extra-scale
screening \cite{gqumond}. Thus neither an additional length scale nor a nonlocal
QUMOND action is claimed here as a new architecture.

Our narrower question can fail: once an isotropic filter changes the amplitude
of the central response, can it tune the isotropic and anisotropic parts
independently? In the linear regime below, the answer is no. The resulting
relation provides a joint consistency test of this candidate, provided both
components can eventually be measured.

\newpage
\section{One static action and its variations}
Let $u$ be an auxiliary Newtonian potential and $\Phi$ the potential felt by
nonrelativistic matter. Fix $a_0>0$ and a real, normalized, isotropic convolution
$S$ on space. Its Fourier multiplier is $\sigma(k)$, with $\sigma(0)=1$.
The filter is self-adjoint under the boundary prescription and commutes
with spatial derivatives. A Gaussian example is
$S=\exp(\xi^2\Delta/2)$, $\sigma(k)=\exp(-\xi^2k^2/2)$.

Use the exact inverse constitutive pair
\begin{equation}
 \mu(y)=1-e^{-y},\qquad s=y\mu(y),\qquad \nu(s)=y/s,\quad y>0.
 \label{eq:pair}
\end{equation}
The inverse is unique since $\dd s/\dd y=1+(y-1)e^{-y}>0$ for $y>0$.
The exponential function itself is not new; it appears explicitly in
earlier Solar-System MOND work \cite{hees}. Define
\begin{align}
 \mathcal G(y)&=y^2+2(1+y)e^{-y}-2, & \mathcal G'(y)&=2s,\\
 q(s^2)&=2sy-\mathcal G(y)-s^2, & q'(s^2)&=\nu(s)-1.
 \label{eq:primitive}
\end{align}
Here $q'$ differentiates its single argument, not $s$. Consider
\begin{equation}
 I=\int\dd t\,\dd^3x\left\{\mathcal L_{m,\mathrm{kin}}-\rho_b\Phi
 -\frac{2\nabla\Phi\cdot\nabla u-|\nabla u|^2
 -a_0^2q\!\left(|\nabla Su|^2/a_0^2\right)}{8\pi G}\right\}.
 \label{eq:action}
\end{equation}
This is a spatially nonlocal, nonrelativistic action, not an Einstein-Hilbert
completion. It specializes Milgrom's functional construction with
$P[u]=|\nabla u|^2/a_0^2+q(|\nabla Su|^2/a_0^2)$ \cite{gqumond}.

Set $f(p)=[\nu(|p|/a_0)-1]p$. Compatible fixed-boundary variations give
\begin{align}
 \frac{\delta I}{\delta\Phi}&=\frac{\Delta u}{4\pi G}-\rho_b=0,\\
 \frac{\delta I}{\delta u}&=\frac{\Delta\Phi-\Delta u-S^*\nabla\cdot f(\nabla Su)}{4\pi G}=0.
 \label{eq:variation}
\end{align}
Consequently,
\begin{equation}
 \boxed{\Delta u=4\pi G\rho_b,\qquad
 \Delta\Phi=4\pi G\rho_b+S^*\nabla\cdot f(\nabla Su).}
 \label{eq:field}
\end{equation}
In particular, $\int f\cdot\nabla S\delta u=-\int\delta u\,S^*\nabla\cdot f$:
the outer filter is required by variation. Matter-coordinate variation gives
acceleration $-\nabla\Phi$; it does not establish a four-dimensional matter
Ward identity. The unsmoothed Newtonian term is retained exactly.

For infinite space these are variational field equations with fixed
asymptotics or background subtraction, not a claim of finite unsubtracted
action. At $S=1$ they reduce to QUMOND. Finite filtering preserves the chosen
constitutive function, \emph{not} the original universal AQUAL equation for
$\Phi$. Even unfiltered QUMOND and AQUAL are not generally identical outside
the appropriate symmetric configurations.

\newpage
\section{A central-tide consistency identity}
Take a uniform background $p_e=\nabla Su_e=a_0s_e\ee$, with $s_e>0$.
The filtered source perturbation must be small relative to $|p_e|$ throughout
the response region. Define
\begin{equation}
 A=\nu(s_e)-1,\quad B=s_e\nu'(s_e),\quad
 Df(p_e)_{ij}=A\delta_{ij}+B e_i e_j.
 \label{eq:AB}
\end{equation}
All following source responses are first order in source amplitude, at fixed
background. For a point-source perturbation of mass $M$, \eqref{eq:field} gives
\begin{equation}
 \widehat{\Phi}_{\ph}(\boldsymbol k)
 =-\frac{4\pi GM}{k^2}|\sigma(k)|^2
 \left[A+B(\kk\cdot\ee)^2\right],\qquad k\ne0,
 \label{eq:fourier}
\end{equation}
where $\Phi_{\ph}$ subtracts the Newtonian source potential and the prescribed
external background. The inverse Fourier measure is $\dd^3k/(2\pi)^3$.
Choose a decaying perturbation at infinity and no additional external tidal
harmonic. Neither the potential constant nor the background is determined
by inverting the $k\ne0$ equation.

\begin{proposition}
Assume isotropic $S$, a spherical source, the nonzero uniform background
and linear regime above, and a convergent central spectral integral. For a
point source put $I_\sigma=\int |\sigma(k)|^2\dd^3k/(2\pi)^3$.
The central phantom-potential Hessian is
\begin{equation}
 H_{ij}=\partial_i\partial_j\Phi_{\ph}(0)
 =\frac{4\pi GM I_\sigma}{15}
 \left[(5A+B)\delta_{ij}+2B e_i e_j\right].
 \label{eq:H}
\end{equation}
With $H_\parallel=e_iH_{ij}e_j$, $H_\perp=(\tr H-H_\parallel)/2$,
$Q_2=H_\perp-H_\parallel$ and $D=\tr H$, one has
\begin{equation}
 \boxed{(15A+5B)Q_2+2BD=0.}
 \label{eq:identity}
\end{equation}
For a finite spherical source, replace $MI_\sigma$ everywhere by
$J_\rho=\int\widehat\rho_b(k)|\sigma(k)|^2\dd^3k/(2\pi)^3$, assuming absolute
convergence. The same identity holds, even if $J_\rho=0$.
\end{proposition}
\begin{proof}
Two derivatives in \eqref{eq:fourier} contribute $-k_i k_j$, canceling the
potential's minus sign. Isotropy separates the radial integral from
\begin{equation}
 \langle n_i n_j\rangle=\frac{\delta_{ij}}3,\qquad
 \langle n_i n_j(n\cdot\ee)^2\rangle
 =\frac{\delta_{ij}+2e_i e_j}{15}.
\end{equation}
These angular moments give \eqref{eq:H}. Its axial difference and trace are
\begin{equation}
 Q_2=-\frac{8\pi GM I_\sigma B}{15},\qquad
 D=4\pi GM I_\sigma\left(A+\frac B3\right).
 \label{eq:QD}
\end{equation}
Substitution proves \eqref{eq:identity}. A spherical finite source changes
only the radial integral, establishing the final statement.
\end{proof}

\newpage
\section{The exponential prediction and its interpretation}
Write $s_e=y(1-e^{-y})$ and $\lambda(y)=1+(y-1)e^{-y}$. Differentiating the
inverse relation, rather than assigning its response, yields
\begin{equation}
 A=\frac{e^{-y}}{1-e^{-y}},\qquad
 B=-\frac{y e^{-y}}{(1-e^{-y})\lambda(y)}.
\end{equation}
The result for the exponential constitutive law is therefore
\begin{equation}
 \boxed{\frac{Q_2}{D}=\frac{2y}{5[3\lambda(y)-y]}}\quad(D\ne0).
 \label{eq:ratio}
\end{equation}
Mass, isotropic filter shape and smoothing length cancel. At $y=1$ the ratio
is $1/5$; as $y\to0^+$ it approaches $2/25$. The latter is generic for a regular
deep-MOND inverse asymptotic $\nu(s)\sim s^{-1/2}$ with its corresponding
derivative asymptotic, not a unique exponential fingerprint. Taking this
limit while retaining external dominance also requires reducing the source
amplitude; the exactly zero-background problem is different.

For reference, the point-source amplitudes are
\begin{equation}
 I_{\mathrm{Gaussian}}=\frac{1}{8\pi^{3/2}\xi^3},\qquad
 I_{\mathrm{Helmholtz}}=\frac{1}{8\pi\xi^3},
\end{equation}
for $\sigma=e^{-\xi^2k^2/2}$ and $\sigma=(1+\xi^2k^2)^{-1}$ respectively.
A general spherical density need not give positive $J_\rho$ for every real
isotropic filter. Nonnegative real-space Gaussian or Helmholtz kernels and
nonnegative density secure positivity for a nonzero source. The tested finite
sources are Gaussian. Positivity is not needed for the identity.

\begin{center}
\begin{tabular}{lll}
\toprule
Point-source feature & Root equation & Computed $y$\\
\midrule
$H_\parallel=0$ & $3y=5\lambda(y)$ & 1.8907779792\\
$D=0$ & $y=3\lambda(y)$ & 3.2602387018\\
$H_\perp=0$ & $y=5\lambda(y)$ & 5.1228414320\\
\bottomrule
\end{tabular}
\end{center}
At the middle root the Hessian is finite; only the quotient in
\eqref{eq:ratio} is undefined. Equation \eqref{eq:identity} is the primary
statement. Between the first and third roots the additional acceleration
stretches parallel to the background and compresses perpendicular to it.
This is not a claim that total gravity is repulsive: the physical tidal
acceleration tensor is $-H$, and the ordinary Newtonian force remains.

Locally the anomalous potential has the expansion
\begin{equation}
 \Phi_{\ph}(x)=\Phi_{\ph}(0)+\frac{D r^2}{6}
 -\frac{Q_2}{2}\left[(\ee\cdot x)^2-\frac{r^2}{3}\right]+o(r^2).
 \label{eq:local}
\end{equation}
Thus $D=4\pi G\rho_{\ph}(0)$ is isotropic central curvature, not total
phantom mass. Our $Q_2$ convention matches the traceless coefficient in
Hees et al., equation (6) \cite{hees}. The extra $D$ term and the weak-source
assumption make this a different prediction from a nonlinear solar
quadrupole calculation; an existing pure-quadrupole fit cannot simply be
read as a fit of both coefficients here.

\newpage
\section{Reproducible checks, not observational evidence}
The companion calculation computes the directional Hessian by numerical
angular integration rather than setting $Q_2/D$ to its expected value.
SymPy separately checks exact angular moments, the inverse derivative,
the division-free identity and the deep limit. Finite differences of the
nonlinear flux check the constitutive Jacobian independently.
\begin{center}
\begin{tabular}{ll}
\toprule
Input & Tested range\\
\midrule
Background $y$ & $0.03,\ 0.3,\ 1,\ 2.5,\ 4,\ 6$\\
Filters & Gaussian, Helmholtz, $1/[1+(\xi k)^4]$\\
Widths $\xi$ & $0.4,\ 1,\ 2.5$ (common numerical units)\\
Gaussian source radii $R$ & $0,\ 0.3,\ 2$ ($R=0$: point source)\\
Directions & $(0,0,1)$ and $(1,2,3)/\sqrt{14}$\\
Full combinations & $6\times3\times3\times3\times2=324$\\
\bottomrule
\end{tabular}
\end{center}
The code sets $G=M=1$ to calculate linear-response coefficients; this is an
amplitude normalization, not a claim that each dimensional realization is
external-field dominated. Finite sources use
$\widehat\rho_b(k)/M=\exp(-R^2k^2/2)$.

The maximum normalized identity residual is $1.10\times10^{-14}$, and the
maximum relative tensor discrepancy is $1.99\times10^{-14}$. The former
divides the absolute left-hand side of \eqref{eq:identity} by
$(|15A+5B|+2|B|)\|H\|_F$; the latter is
$\|H_{\rm num}-H_{\rm analytic}\|_F/\|H_{\rm analytic}\|_F$.
These are floating-point agreement measures, not empirical error bars or
certified continuum error bounds.

Gaussian and Helmholtz radial integrals are checked against exact
normalizations. Increasing the polar quadrature from 24 to 40 nodes and
rotating the field axis checks the angular implementation. An intentionally
anisotropic spectral weight $1+0.8(\kk\cdot\ee)^2$ gives an identity residual
of $0.06672$ at $y=1$. This negative control illustrates why isotropy is a
load-bearing hypothesis.

The current rerun passes six tidal tests, ten parent-action tests, six
existing exact-inverse-kernel tests and two elliptic-action regression tests:
24 distinct tests, all with exit status zero. The parent checks include a
finite-difference first variation. They do not collectively certify a
relativistic field theory. The numerical environment is Python 3.9.6,
NumPy 1.26.2, SciPy 1.11.4 and SymPy 1.14.0. The companion manifest records
runtime, repository state and input/output checksums; the adjacent reproduction
note gives commands. No astronomical observations were fitted or analyzed.

\section{What a falsification would require}
The first practical calculation is a fully nonlinear solution of
\eqref{eq:field} with a known external source, testing how the relation changes
as the filtered internal field ceases to be small. The host system must also
calibrate $s_e$: $y$ is an inverse-kernel parameter, not automatically the
measured physical external acceleration divided by $a_0$.
Milgrom's external-boundary analysis distinguishes those fields in a general
host \cite{qumond}. A meaningful data test then needs joint identifiability
of $D$ and $Q_2$, including ordinary tides, source shape and nuisance
parameters. Increasing $\xi$ can suppress both amplitudes together; the ratio
alone does not exclude that escape.

\newpage
\section{Scope, priority and the unfinished relativistic problem}
The proposition is an analytic consequence of the specified static action,
its constant-background linearization and isotropic angular averages. It is
not inferred from 324 numerical successes. Those checks do not test nonlinear
source corrections, nonspherical systems or varying external fields. The
exactly zero-background limit is singular in the constitutive tangent; finite
central curvature at nonzero background does not settle strong coupling or
loss of rank at zero field.

There is no physical spacetime metric in \eqref{eq:action}. This paper derives
neither $\Phi=\Psi$ nor lensing, PPN parameters, a two-tensor Dirac count,
covariant matter conservation, gravitational-wave speed, perturbative
stability or FLRW evolution. Elliptic potentials are not proof that a
covariant completion has no hidden mode or instantaneous physical channel.
Both $a_0$ and $\xi$ are inputs; no cosmological acceleration-scale relation
is derived. The exact exponential constitutive choice is retained, but the
stronger original AQUAL target is not fulfilled by this finite-filter action.

The explicit functional architecture has prior art \cite{gqumond}. The central
filtered identity was not identified in the inspected source passages or the
bounded repository/web searches documented with this draft. That is not
evidence of global priority. The defensible contribution is an explicitly
derived and reproducibly checked central-response corollary for this
specialization. No claim of a first discovery is made.

\section{Conclusion}
The external-field anisotropy emphasized by Banik and Zhao survives isotropic
smoothing in a constrained form: changing the filter amplitude cannot
independently change the two central tidal coefficients in the stated linear
regime. This is a concrete, conditional prediction from one action, not a
fitted pair of unrelated signals. Its scientific value will be decided by
nonlinear corrections and joint observational identifiability. The static
identity is established under its assumptions; the full theory remains open.

\begin{thebibliography}{9}
\bibitem{banik}
I. Banik and H. Zhao,
\emph{The External Field Dominated Solution In QUMOND \& AQUAL:
Application To Tidal Streams},
\href{https://arxiv.org/pdf/1509.08457v2}{arXiv:1509.08457v2},
4 May 2018 revision; equations (22)--(25), (36).
\bibitem{qumond}
M. Milgrom, \emph{Quasi-linear formulation of MOND},
\href{https://arxiv.org/pdf/0911.5464v2}{arXiv:0911.5464v2},
2 March 2010 revision; equations (3)--(6) and section 5.
\bibitem{gqumond}
M. Milgrom, \emph{Generalizations of quasilinear MOND (QUMOND)},
\href{https://arxiv.org/pdf/2305.01589v2}{arXiv:2305.01589v2},
4 October 2023 revision; section II.B and equations (6)--(7).
\bibitem{hees}
A. Hees, W. M. Folkner, R. A. Jacobson and R. S. Park,
\emph{Constraints on MOND theory from radio tracking data of the Cassini
spacecraft}, \href{https://arxiv.org/pdf/1402.6950v2}{arXiv:1402.6950v2},
29 April 2014 revision; equations (5b) and (6).
\end{thebibliography}
\small\textit{Research integrity note.} Drafting, symbolic work and verification
were AI-assisted. This is not a journal acceptance, an endorsement by the
cited authors, or a claim to have settled observational disputes. Human
author review is required before submission or public dissemination.
\end{document}
